% ============================================================
% 第五讲 · Part 02：需求概念与最优值的存在性
% ============================================================

\section{需求概念与上确界：Marshallian 需求、概念家族与存在性}
\label{sec:demand-concepts}

效用最大化问题的解——\textbf{需求}——是消费者理论的核心输出。
本节首先厘清\textbf{需求对应}与\textbf{需求函数}的区别，
然后系统梳理围绕 UMP 构建的整个\textbf{需求概念家族}。

\subsection{Marshallian 需求：对应与函数}
\label{subsec:correspondence-vs-function}

\begin{definition}[Marshallian（Walrasian）需求对应]\label{def:marshallian-demand}
设 $u:\R^L_+\to\R$ 为效用函数，$p\gg 0$ 为价格向量，$w\ge 0$ 为财富。
\textbf{Marshallian 需求对应}定义为 UMP 的最优解集：
\[
  X(p,w)=\arg\max\bigl\{\,u(x)\;:\;x\in\R^L_+,\;p\cdot x\le w\,\bigr\}.
\]
当 $X(p,w)$ 对给定的 $(p,w)$ 恰好只含一个元素时，
记 $x(p,w)$ 为其中唯一的消费束，称为\textbf{Marshallian 需求函数}。
\end{definition}

$X(p,w)$ 之所以称为\textbf{对应}（correspondence）而非\textbf{函数}（function），
是因为从 $(p,w)$ 到消费束的映射可能是\emph{多值}的：
同一个价格—财富组合下，可能存在多个互不相同的最优消费束。
何时出现多值？何时退化为单值？
答案是：取决于偏好的\textbf{凸性}。

\begin{itemize}
\item \textbf{偏好凸但非严格凸}（如完全替代偏好）：
  无差异曲线可能含有\emph{线性段}。
  当预算线的斜率恰好等于线性段的斜率时，
  线性段上\emph{每一点}都与预算线"相切"，
  整个线段上的消费束同时达到最大效用，
  此时 $X(p,w)$ 是一段线段（多值）。
\item \textbf{偏好严格凸}（如 Cobb--Douglas 偏好）：
  无差异曲线处处严格弯曲，不含线性段。
  每条无差异曲线与预算线至多只有一个接触点，
  因而 $X(p,w)$ 退化为单点集 $\{x(p,w)\}$，
  需求成为\emph{函数}。
\end{itemize}

图~\ref{fig:convex-vs-strict} 直观展示两种情形的差异。

\begin{figure}[H]
\centering
\begin{tikzpicture}[scale=0.78, line cap=round, line join=round]
  % ===== 左图：凸偏好（非严格凸），需求对应为一段线段 =====
  \begin{scope}
    % 坐标轴
    \draw[axisline] (0,0) -- (6.2,0) node[below, axislabel] {$x_1$};
    \draw[axisline] (0,0) -- (0,5.8) node[left, axislabel] {$x_2$};
    % 预算线：从 (0,4.5) 到 (5.5,0)，斜率 -4.5/5.5 ≈ -0.818
    \draw[thick] (0,4.5) -- (5.5,0);
    \node[figlabel, anchor=south west] at (0.1,4.55) {$w/p_2$};
    \node[figlabel, anchor=north east] at (5.45,-0.05) {$w/p_1$};
    % 无差异曲线：上弯部分 + 线性段（与预算线重合）+ 下弯部分
    \draw[curveA] (0.4, 5.3) .. controls (0.6, 4.2) and (1.0, 3.1) .. (1.6, 2.55);
    \draw[curveA, line width=1.8pt] (1.6, 2.55) -- (3.8, 0.75);
    \draw[curveA] (3.8, 0.75) .. controls (4.3, 0.44) and (4.8, 0.32) .. (5.6, 0.22);
    % 最优线段端点标记
    \fill (1.6, 2.55) circle (2.2pt);
    \fill (3.8, 0.75) circle (2.2pt);
    \node[figlabel, anchor=south east] at (1.5, 2.7) {$x^a$};
    \node[figlabel, anchor=north west] at (3.9, 0.6) {$x^b$};
    % 标注：整段最优
    \node[align=center, font=\footnotesize] at (2.7, 3.6)
      {凸（非严格凸）偏好：\\ $X(p,w)=[x^a,x^b]$ 为线段};
    % 一条更低的无差异曲线（虚线）
    \draw[curveB] (0.5, 3.6) .. controls (1.0, 2.2) and (2.0, 1.3) .. (3.2, 0.85);
    \draw[curveB] (3.2, 0.85) -- (4.6, 0.30);
    \draw[curveB] (4.6, 0.30) .. controls (4.9, 0.20) and (5.2, 0.15) .. (5.7, 0.10);
    \node[curveB] at (5.2, 0.75) {};
  \end{scope}
  % ===== 右图：严格凸偏好，唯一最优点 =====
  \begin{scope}[xshift=7.4cm]
    % 坐标轴
    \draw[axisline] (0,0) -- (6.2,0) node[below, axislabel] {$x_1$};
    \draw[axisline] (0,0) -- (0,5.8) node[left, axislabel] {$x_2$};
    % 预算线：从 (0,4.5) 到 (5.5,0)
    \draw[thick] (0,4.5) -- (5.5,0);
    \node[figlabel, anchor=south west] at (0.1,4.55) {$w/p_2$};
    \node[figlabel, anchor=north east] at (5.45,-0.05) {$w/p_1$};
    % 严格凸无差异曲线：x₂ = k/x₁ 型，与预算线相切
    % 切点在 (2.5, 2.45)（预算线上），k = 2.5×2.45 = 6.125
    \draw[curveA] plot [smooth, tension=0.8]
      coordinates {(1.2, 5.1) (1.5, 4.08) (2.0, 3.06) (2.5, 2.45) (3.0, 2.04) (3.5, 1.75) (4.2, 1.46) (5.2, 1.18)};
    % 最优点标记
    \fill (2.5, 2.45) circle (2.2pt);
    \node[figlabel, anchor=north west] at (2.65, 2.30) {$x^* = x(p,w)$};
    % 标注
    \node[align=center, font=\footnotesize] at (2.7, 3.9)
      {严格凸偏好：\\ $X(p,w)=\{x^*\}$ 单点};
    % 一条更低的无差异曲线（虚线）：k = 3.5
    \draw[curveB] plot [smooth, tension=0.8]
      coordinates {(1.0, 3.5) (1.4, 2.5) (2.0, 1.75) (2.8, 1.25) (3.5, 1.0) (4.5, 0.78) (5.5, 0.64)};
  \end{scope}
\end{tikzpicture}
\caption{偏好凸性与需求对应的单值性：左图为凸但非严格凸偏好，
  需求对应为一段线段；右图为严格凸偏好，需求退化为唯一的切点。}
\label{fig:convex-vs-strict}
\end{figure}

\figexplan{
左图：无差异曲线（实线）在中间段为直线（线性段），且线性段恰好落在预算线上。
此时线性段两端点 $x^a$ 与 $x^b$ 之间（含端点）的每个消费束，
其所在的无差异曲线都"同时"与预算线重合，效用值相同且均达到最优，
因此 Marshallian 需求对应 $X(p,w)$ 是整个线段 $[x^a,x^b]$。
虚线为一条更低的无差异曲线，其上的消费束同样可负担但未达最优。
右图：无差异曲线处处严格弯曲（不含线性段），
与预算线只在唯一切点 $x^*$ 处接触。
该无差异曲线是预算集内所能达到的最高无差异曲线，
切点 $x^*$ 是唯一的最优消费束，故需求对应退化为需求函数 $x(p,w)$。
虚线为一条更低的无差异曲线，不与预算线相切。}

\subsection{需求概念家族一览}
\label{subsec:demand-family}

围绕 UMP，微观经济学构建了一族紧密关联的概念。
表~\ref{tab:demand-family} 给出完整的概览。

\begin{definition}[间接效用函数、Hicksian 需求与支出函数]\label{def:indirect-hicksian}
设 $u:\R^L_+\to\R$ 连续，$p\gg 0$，$w\ge 0$，目标效用 $\bar u$ 满足可达性。
\begin{enumerate}
\item \textbf{间接效用函数}（indirect utility function）：
  Marshallian 需求所对应的效用水平
  \[
    v(p,w)=\max\bigl\{\,u(x)\;:\;x\in\R^L_+,\;p\cdot x\le w\,\bigr\}
      = u\bigl(x(p,w)\bigr).
  \]
\item \textbf{Hicksian（补偿）需求}（compensated demand）：
  在\emph{达到目标效用}的前提下，最小化支出
  \[
    h(p,\bar u)=\arg\min\bigl\{\,p\cdot x\;:\;x\in\R^L_+,\;u(x)\ge\bar u\,\bigr\}.
  \]
\item \textbf{支出函数}（expenditure function）：
  上述最小化问题的最优值
  \[
    e(p,\bar u)=\min\bigl\{\,p\cdot x\;:\;x\in\R^L_+,\;u(x)\ge\bar u\,\bigr\}
      = p\cdot h(p,\bar u).
  \]
\end{enumerate}
\end{definition}

\begin{table}[H]
\centering
\small
\begin{tabular}{@{}llll@{}}
\toprule
\textbf{概念} & \textbf{数学定义} & \textbf{固定什么} & \textbf{回答什么问题} \\
\midrule
Marshallian 需求 & $\max u(x)$ s.t.\ $p\cdot x\le w$ & 价格 $p$、财富 $w$ & 给定预算，买什么？ \\
$x(p,w)$ & 的最优解 & & \\
\addlinespace[3pt]
间接效用 $v(p,w)$ & 同上问题的最优值 & 价格 $p$、财富 $w$ & 给定预算，最高效用是多少？ \\
\addlinespace[3pt]
Hicksian 需求 $h(p,\bar u)$ & \multicolumn{1}{l}{$\min p\cdot x$ s.t.\ $u(x)\ge\bar u$} & 价格 $p$、目标效用 $\bar u$ & 达到目标效用，买什么？ \\
 & 的最优解 & & \\
\addlinespace[3pt]
支出函数 $e(p,\bar u)$ & 同上问题的最优值 & 价格 $p$、目标效用 $\bar u$ & 达到目标效用，最少花多少钱？ \\
\bottomrule
\end{tabular}
\caption{需求概念家族一览}
\label{tab:demand-family}
\end{table}

这一家族的核心联系由如下\textbf{对偶关系}给出：
\begin{equation}\label{eq:duality-relations}
  h(p,\bar u) = x\bigl(p,\,e(p,\bar u)\bigr),
  \qquad
  e\bigl(p,\,v(p,w)\bigr) = w.
\end{equation}
第一式含义：最小化支出得到的最优消费束，
恰好等于"给消费者恰好够用的财富 $e(p,\bar u)$ 让其自己最大化效用"的结果；
第二式含义：最大化效用之后再问"达到 $v(p,w)$ 至少要花多少钱"，
答案就是 $w$。
这两条对偶关系将在后续章节（支出最小化问题）中系统展开，
并构成 Slutsky 分解的理论基础：
\textbf{价格变化的总效应}可以拆分为
\emph{替代效应}（Hicksian 需求的变化，效用不变）与
\emph{财富效应}（由实际购买力变化引起）。
没有 Hicksian 需求与支出函数这对工具，
Slutsky 分解便无从谈起。

\begin{remark}[为什么需要两套概念？]
Marshallian 需求回答"给定预算怎么花"；
Hicksian 需求回答"给定效用怎么最省"。
前者是消费者真实面对的问题，与数据直接对应；
后者是理论分析的理想化问题，结构更简洁。
两者通过 \eqref{eq:duality-relations} 的对偶关系紧密联结：
$v$ 是 $e$ 的"反函数"（在适当的意义下），
$h$ 可由 $e$ 经 Shephard 引理直接恢复，
正如 $x$ 可由 $v$ 经 Roy 恒等式恢复。
对偶结构的完整刻画留待后续章节。
\end{remark}
\label{sec:sup-max}

\subsection{sup 与 max 的严格辨析}
\label{subsec:sup-vs-max}

设 $S\subseteq\R$ 为一实数集，$f$ 为定义在某集合 $D$ 上的实值函数。

\begin{definition}[上确界与最大值]\label{def:sup-max}
\begin{enumerate}
\item $M\in\R$ 是 $S$ 的\textbf{最大值}，若 (i) $M\in S$ 且 (ii) $s\le M$ 对一切 $s\in S$ 成立。
\item $\bar s\in\R$ 是 $S$ 的\textbf{上确界}（supremum），若 (i) $s\le\bar s$ 对一切 $s\in S$（上界），
  且 (ii) 对任意 $\bar s'<\bar s$，存在 $s\in S$ 使 $s>\bar s'$（最小上界）。
\end{enumerate}
最大值必定是上确界；但上确界可能不属于 $S$，此时最大值不存在。
\end{definition}

\begin{example}[sup 存在而 max 不存在]\label{ex:sup-not-max}
\leavevmode
\begin{enumerate}
\item 开区间 $S=(0,1)$：$\sup S=1\notin S$，最大值不存在。
\item 定义在开集 $D=(0,1)$ 上的函数 $f(x)=x$：
  $\sup_{x\in D} f(x)=1$，但 $f$ 在 $D$ 上没有最大值
  （因为 $x=1\notin D$，而 $f(x)$ 可以无限接近但永远不等于 $1$）。
\end{enumerate}
在 UMP 中，预算集 $B(p,w)$ 是\emph{闭}的
（它包含边界 $\{x:p\cdot x=w\}$），
这保证了"最优解不会被遗漏在可行集之外"。
\end{example}

上确界与最大值何时一致？
这正是 \textbf{Weierstrass 定理}回答的问题：
在紧集上，连续函数的最大值必然达到。

\subsection{Weierstrass 定理：完整证明}
\label{subsec:weierstrass}

\begin{theorem}[Weierstrass 最大值定理]\label{thm:weierstrass}
设 $K\subseteq\R^L$ 为非空紧集，$f:K\to\R$ 为连续函数。
则 $f$ 在 $K$ 上达到最大值与最小值：
存在 $x^*,x_*\in K$ 使 $f(x^*)\ge f(x)\ge f(x_*)$ 对一切 $x\in K$ 成立。
\end{theorem}

\begin{proof}
只证明最大值的情形（最小值通过对 $-f$ 取最大值得到）。

\textbf{第一步：$f$ 在 $K$ 上有界。}
反设 $f$ 在 $K$ 上无上界。
则对每个 $n\in\N$，存在 $x^n\in K$ 使 $f(x^n)>n$。
由于 $K$ 是紧集（$\R^L$ 中紧 $\Leftrightarrow$ 闭且有界 $\Leftrightarrow$ 序列紧），
序列 $\{x^n\}$ 存在收敛子列 $\{x^{n_k}\}$，其极限 $x^*=\lim_k x^{n_k}\in K$。
由 $f$ 连续，$f(x^{n_k})\to f(x^*)$。
但 $f(x^{n_k})>n_k\to\infty$，与 $f(x^*)<\infty$ 矛盾。
故 $f$ 在 $K$ 上有上界，从而 $M=\sup_{x\in K}f(x)<\infty$ 存在。

\textbf{第二步：上确界可以达到。}
由上确界的定义（最小上界性质），
对每个 $n\in\N$ 存在 $x^n\in K$ 使 $f(x^n)>M-\frac{1}{n}$。
再次由 $K$ 的序列紧性，
存在收敛子列 $\{x^{n_k}\}$，其极限 $x^*\in K$。
由 $f$ 连续，
\[
  f(x^*)=\lim_{k\to\infty}f(x^{n_k}).
\]
另一方面，$M-\frac{1}{n_k}<f(x^{n_k})\le M$ 对一切 $k$ 成立，
令 $k\to\infty$，由夹逼定理得
\[
  f(x^*)=M=\sup_{x\in K}f(x).
\]
又 $x^*\in K$，故 $M\in f(K)$，
即上确界作为函数值被 $x^*$ 达到，$x^*$ 即为最大点。
\end{proof}

\begin{remark}[证明结构解读]
证明的核心是\textbf{序列紧性}（sequential compactness）：
紧集中的任意序列必有收敛于该集合内的子列。
第一步排除了目标值发散到无穷的可能（保证 sup 有限），
第二步将"无限逼近"转化为"确实达到"。
两步共同的关键前提是：$K$ 是\textbf{闭}的（子列极限不会跑出可行集）
且\textbf{有界}（序列不会跑向无穷远）。
这解释了为何在 UMP 中，
预算集 $B(p,w)$ 的闭性与有界性缺一不可：
闭性要求预算约束包含等号（$\le w$ 而非 $<w$），
有界性要求所有价格为正（$p\gg 0$）。
\end{remark}

\subsection{需求对应的三条基本定理}
\label{subsec:demand-three-theorems}

将 Weierstrass 定理应用于预算集，
我们立即得到需求理论的三条基本定理。

\begin{theorem}[需求对应非空]\label{thm:demand-nonempty}
设 $u:\R^L_+\to\R$ 连续，$p\gg 0$，$w\ge 0$。
则 Marshallian 需求对应 $X(p,w)$ 非空（且为紧集）。
\end{theorem}

\begin{proof}
预算集
\[
  B(p,w)=\bigl\{x\in\R^L_+:p\cdot x\le w\bigr\}
\]
是紧集：
\emph{闭}，因为它是不等式 $p\cdot x-w\le 0$ 的零点下方集
（连续函数的闭集的原像）与非负象限之交；
\emph{有界}，因为 $p_i>0$ 且 $x_i\ge 0$ 对每个 $i$ 成立，
故 $x_i\le w/p_i$。
又 $0\in B(p,w)$，故 $B(p,w)$ 非空。
由 Weierstrass 定理（定理~\ref{thm:weierstrass}），
连续函数 $u$ 在非空紧集 $B(p,w)$ 上达到最大值，
即 $X(p,w)=\arg\max_{B(p,w)}u\neq\varnothing$。

$X(p,w)$ 的紧性：$X(p,w)$ 是闭集
（它是 $\{x\in B(p,w):u(x)=v(p,w)\}$，
即连续函数 $u$ 的水平集与闭集 $B(p,w)$ 之交，故为闭集）
且含于有界集 $B(p,w)$，故 $X(p,w)$ 为紧集。
\end{proof}

\begin{theorem}[凸偏好下需求对应为凸集]\label{thm:demand-convex}
设 $u:\R^L_+\to\R$ 连续且拟凹（代表凸偏好），
$p\gg 0$，$w\ge 0$。
则 $X(p,w)$ 是凸集。
\end{theorem}

\begin{proof}
设 $x^a, x^b\in X(p,w)$，$t\in[0,1]$。
记 $x^t=tx^a+(1-t)x^b$。

\textbf{$x^t$ 可行}：由 $p\cdot x^a\le w$、$p\cdot x^b\le w$ 与 $p$ 的线性性，
\[
  p\cdot x^t = t\,p\cdot x^a + (1-t)\,p\cdot x^b \le t\,w+(1-t)\,w = w,
\]
且 $x^t\in\R^L_+$（非负象限是凸集），故 $x^t\in B(p,w)$。

\textbf{$x^t$ 最优}：由 $u$ 的拟凹性，
\[
  u(x^t) \ge \min\bigl\{u(x^a),\,u(x^b)\bigr\}.
\]
而 $x^a, x^b$ 均为最优解，故 $u(x^a)=u(x^b)=v(p,w)$，其中
\[
  v(p,w)=\max_{x\in B(p,w)}u(x)
\]
为间接效用函数。
于是 $u(x^t)\ge v(p,w)$。
但 $v(p,w)$ 是 $u$ 在 $B(p,w)$ 上的最大值，
故 $u(x^t)\le v(p,w)$。
综合得 $u(x^t)=v(p,w)$，即 $x^t\in X(p,w)$。
\end{proof}

\begin{theorem}[严格凸偏好下需求唯一]\label{thm:demand-unique}
设 $u:\R^L_+\to\R$ 连续且严格拟凹（代表严格凸偏好），
$p\gg 0$，$w\ge 0$。
则 $X(p,w)$ 至多含有一个元素，
即 Marshallian 需求退化为单值的需求函数 $x(p,w)$。
\end{theorem}

\begin{proof}
反设 $X(p,w)$ 含有两个不同元素 $x^a\neq x^b$。
由定理~\ref{thm:demand-nonempty}，$X(p,w)\neq\varnothing$，
故 $v(p,w)$ 有限且 $u(x^a)=u(x^b)=v(p,w)$。

取 $t\in(0,1)$，记 $x^t=tx^a+(1-t)x^b$。
由定理~\ref{thm:demand-convex} 的证明（第一步），
$x^t\in B(p,w)$。
由 $u$ 的\emph{严格}拟凹性，
\[
  u(x^t) > \min\bigl\{u(x^a),\,u(x^b)\bigr\} = v(p,w).
\]
但 $v(p,w)$ 是 $u$ 在 $B(p,w)$ 上的最大值，
$u(x^t)>v(p,w)$ 与此矛盾。
故 $X(p,w)$ 不可能含有两个不同元素。
\end{proof}

\subsection{Walras 定律：预算约束取等}
\label{subsec:walras-law}

在局部非饱和偏好下，消费者永远不会"留钱不花"：
最优消费束必定恰好花完财富。
这一性质称为 Walras 定律，
它是需求理论中最基础的可检验含义之一。

\begin{theorem}[Walras 定律]\label{thm:walras}
设偏好局部非饱和（对任意 $x$ 及任意 $\varepsilon>0$，
存在 $y\in\R^L_+$ 使 $\|y-x\|<\varepsilon$ 且 $y\succ x$），
$u$ 为其连续表示，$p\gg 0$，$w>0$。
则对任意 $x^*\in X(p,w)$，有
\[
  p\cdot x^* = w.
\]
\end{theorem}

\begin{proof}
由 $x^*\in X(p,w)$，显然 $p\cdot x^*\le w$。
反设 $p\cdot x^*<w$，即 $x^*$ 位于预算集的\emph{内部}
（严格满足不等式）。
令
\[
  \varepsilon = \frac{w - p\cdot x^*}{2\|p\|} > 0.
\]
由局部非饱和，存在 $y\in\R^L_+$ 使 $\|y-x^*\|<\varepsilon$ 且 $u(y)>u(x^*)$。

考察 $y$ 的可负担性。由 Cauchy--Schwarz 不等式，
\[
  p\cdot y = p\cdot x^* + p\cdot(y-x^*)
  \le p\cdot x^* + \|p\|\,\|y-x^*\|
  < p\cdot x^* + \|p\|\,\varepsilon
  = p\cdot x^* + \frac{w-p\cdot x^*}{2}
  < w.
\]
故 $y\in B(p,w)$ 且 $u(y)>u(x^*)$。
这与 $x^*$ 是 $B(p,w)$ 上的最大点矛盾。
故必有 $p\cdot x^*=w$。
\end{proof}

图~\ref{fig:walras-interior} 给出几何直观：
若消费束位于预算集内部（$p\cdot x<w$），
则存在包含它的一个小邻域仍然在预算集内；
由局部非饱和，邻域内总存在严格更好的消费束；
因此内部点不可能是最优的——最优点必须落在预算集的边界上，
即恰好花完财富。

\begin{figure}[H]
\centering
\begin{tikzpicture}[scale=0.82, line cap=round, line join=round]
  % 坐标轴
  \draw[axisline] (0,0) -- (7.2,0) node[below, axislabel] {$x_1$};
  \draw[axisline] (0,0) -- (0,6.2) node[left, axislabel] {$x_2$};
  % 预算线
  \draw[thick] (0,5) -- (6,0);
  \node[figlabel, anchor=south west] at (0.1,5.05) {$w/p_2$};
  \node[figlabel, anchor=north east] at (5.95,-0.05) {$w/p_1$};
  % 内点 x（非最优候选）
  \fill (2.5, 2.0) circle (2.2pt);
  \node[figlabel, anchor=south east] at (2.3, 2.2) {$x$};
  % 邻域（小圆）
  \draw[dashed] (2.5, 2.0) circle (1.2);
  \node[figlabel, anchor=north west] at (3.7, 3.1) {\small $\varepsilon$-邻域仍在 $B(p,w)$ 内};
  % 邻域内的更优点 y
  \fill (3.0, 2.5) circle (2.2pt);
  \node[figlabel, anchor=south west] at (3.2, 2.7) {$y\succ x$};
  % 箭头指向更优点
  \draw[arrblue, thick] (2.65, 2.15) -- (2.85, 2.35);
  % 预算边界上的最优解
  \fill (4.0, 1.67) circle (2.2pt);
  \node[figlabel, anchor=west] at (4.2, 1.67) {$x^*\in X(p,w)$：必须在边界上};
  % 过 x 的无差异曲线（示意，虚线）
  \draw[curveB] (1.2, 4.2) .. controls (2.0, 3.0) and (3.0, 2.6) .. (4.2, 2.4);
  \draw[curveB] (4.2, 2.4) .. controls (5.0, 2.2) and (5.8, 2.1) .. (6.8, 2.0);
\end{tikzpicture}
\caption{Walras 定律的几何直观：内部点不可能是最优解。
  若 $p\cdot x<w$，则 $x$ 的一个小邻域整体落在预算集内，
  局部非饱和保证邻域内存在严格更好的消费束 $y$，
  因而 $x$ 不可能最优。最优解必须落在预算线上，
  即 $p\cdot x^*=w$。}
\label{fig:walras-interior}
\end{figure}

\figexplan{
预算线（实线）将可行区域（左下方）与不可行区域分开。
候选消费束 $x$（左点）严格位于预算集内部，即 $p\cdot x<w$。
以 $x$ 为中心、适当小的半径画一个小圆（虚线），
由 $p\cdot x<w$ 且不等式是严格的，可将半径选得足够小使整个圆仍在预算集内。
由局部非饱和假设，圆内必存在严格优于 $x$ 的消费束 $y$（右点），
且 $y$ 同样可负担。
因此 $x$ 被支配，不可能是最优解。
反观位于预算线上的点 $x^*$（右侧圆点），
预算线之外不存在任何可负担的消费束，
故它可能成为最优解。
这从几何上解释了 Walras 定律：
在局部非饱和偏好下，消费者将财富恰好用尽。
}

\begin{remark}[Walras 定律的含义与应用]
Walras 定律 $p\cdot x^*=w$ 有两个重要推论。
其一，\textbf{需求的零次齐次性}：
当价格与财富同比例变化 $(p,w)\to(\lambda p,\lambda w)$（$\lambda>0$）时，
预算集 $B(p,w)$ 不变，因此 $X(p,w)=X(\lambda p,\lambda w)$。
消费者不受"货币幻觉"影响，只关心相对价格与实际购买力。
其二，\textbf{加总与市场层面}：
将所有消费者的 Walras 定律加总，
可得市场总超额需求的价值恒为零。
这一性质在一般均衡理论中用于证明：
若 $L-1$ 个市场出清，则第 $L$ 个市场自动出清
（Walras 定律在一般均衡中的版本）。
\end{remark}

\begin{remark}[对偶预告]
本节建立的三条基本定理（非空性、凸值性、单值性）与 Walras 定律
构成了需求理论的公理化骨架。
后续章节将在此骨架上叠加微积分结构：
对偶关系 \eqref{eq:duality-relations} 将给出
Hicksian 与 Marshallian 需求之间的精确换算（Slutsky 方程），
间接效用函数 $v$ 的可微性将产生 Roy 恒等式，
支出函数 $e$ 的可微性将产生 Shephard 引理。
这些结果共同构成现代需求理论的完整体系。
\end{remark}
